\documentclass[../../main.tex]{subfiles}% 注意这里的文件路径不能用 ./main.tex ,否则用latexmk编译子文件会报错
\graphicspath{{\subfix{./image/}}{\subfix{../../image/}}} % 指定图片引用路径,后续可以直接使用图片文件名
% 如果图片存放在上级目录,则使用 ../../image/ 作为备用路径

% 例如:
% \begin{figure}[H]
% \centering
% \includegraphics[width=0.3\textwidth]{图.png}
% \caption{}
% \label{figure:图}
% \end{figure}
% 注意:上述\label{}一定要放在\caption{}之后,否则引用图片序号会只会显示??.

\begin{document}

\section{正规族}

\begin{definition}
两个域 \(D_1\) 和 \(D_2\) 称为是\textbf{全纯等价的},如果存在单叶的全纯函数 \(f: D_1 \to \bm{C}\),使得 \(f(D_1)=D_2\),即 \(f\) 一一地把 \(D_1\) 映为 \(D_2\). 
\end{definition}
\begin{remark}
显然,全纯等价是一个等价关系. 
\end{remark}

\begin{definition}
设 \(\mathscr{F}\) 是域 \(D\) 上的一个函数族,如果它的任意序列 \(\{f_n\} \subseteq \mathscr{F}\) 中一定包含一个在 \(D\) 上内闭一致收敛的子列 \(\{f_{n_k}\}\),就称 \(\mathscr{F}\) 是 \(D\) 上的一个\textbf{正规族}.
\end{definition}

\begin{definition}
设 \(\mathscr{F}\) 是域 \(D\) 上的一个函数族,如果存在常数 \(M\),使得对任意的 \(f \in \mathscr{F}\) 及 \(z \in D\),均有 \(|f(z)| \leqslant M\),就说 \(\mathscr{F}\) 在 \(D\) 上是\textbf{一致有界的}.

如果对任意紧集 \(K \Subset D\),存在与 \(K\) 有关的常数 \(M(K)\),使得对任意的 \(f \in \mathscr{F}\) 及 \(z \in K\),有 \(|f(z)| \leqslant M(K)\),就说 \(\mathscr{F}\) 在 \(D\) 上是\textbf{内闭一致有界的}. 
\end{definition}
\begin{remark}
显然,在 \(D\) 上一致有界的函数族一定是内闭一致有界的,反之不然.
\end{remark}

\begin{definition}
设 \(\mathscr{F}\) 是域 \(D\) 上的一个函数族,如果对任意 \(\varepsilon > 0\),存在 \(\delta > 0\),当 \(z_1, z_2 \in D\),且 \(|z_1 - z_2| < \delta\) 时,有 \(|f(z_1) - f(z_2)| < \varepsilon\) 对每个 \(f \in \mathscr{F}\) 都成立,就说 \(\mathscr{F}\) 在 \(D\) 上是\textbf{等度连续的}.
\end{definition}

\begin{theorem}\label{thm:arzela_ascoli}
设 \(K\) 是 \(\bm{C}\) 中的紧集,\(\{f_n\}\) 是在 \(K\) 上一致有界且等度连续的函数列,那么 \(\{f_n\}\) 必有子列在 \(K\) 上一致收敛.
\end{theorem}
\begin{proof}
令 \(A \subseteq K\) 是可数集,且在 \(K\) 中稠密. 记
\begin{align*}
A = \{\zeta_1, \zeta_2, \cdots\}.
\end{align*}
由于 \(\{f_n(\zeta_1)\}\) 是有界数列,根据\hyperref[theorem:Bolzano-Weierstrass定理]{Bolzano-Weierstrass定理},从中可以取出收敛的子列,即有
\begin{align*}
f_{n_1}^{(1)}(z), f_{n_2}^{(1)}(z), \cdots,
\end{align*}
在 \(z = \zeta_1\) 处收敛;因为 \(\{f_{n_j}^{(1)}(\zeta_2)\}\) 有界,故有
\begin{align*}
f_{n_1}^{(2)}(z), f_{n_2}^{(2)}(z), \cdots,
\end{align*}
在 \(z = \zeta_2\) 处收敛. 继续这样做下去,可得一串函数列:
\begin{align*}
f_{n_1}^{(1)}(z), f_{n_2}^{(1)}(z), \cdots, f_{n_k}^{(1)}(z), \cdots; \\
f_{n_1}^{(2)}(z), f_{n_2}^{(2)}(z), \cdots, f_{n_k}^{(2)}(z), \cdots; \\
\cdots; \\
f_{n_1}^{(s)}(z), f_{n_2}^{(s)}(z), \cdots, f_{n_k}^{(s)}(z), \cdots;
\end{align*}
其中,后一序列是前一序列的子列,第 \(s\) 个序列在 \(\zeta_1, \zeta_2, \cdots, \zeta_s\) 处收敛. 现在取对角线序列
\begin{align}
f_{n_1}^{(1)}(z), \cdots, f_{n_2}^{(2)}(z), \cdots, f_{n_k}^{(k)}(z), \cdots, \label{eq:diag_seq}
\end{align}
它在 \(A\) 中的每一点处都收敛. 为符号简单起见,我们把函数列 \eqref{eq:diag_seq} 记为 \(\{f_{n_k}\}\),它是 \(\{f_n\}\) 的一个子列,我们证明它在 \(K\) 上一致收敛. 因为 \(\{f_n\}\) 在 \(K\) 上等度连续,故对任意 \(\varepsilon > 0\),存在 \(\delta > 0\),当 \(z_1, z_2 \in K\),且 \(|z_1 - z_2| < \delta\) 时,有 \(|f_n(z_1) - f_n(z_2)| < \varepsilon\) 对 \(n=1, 2, \cdots\) 都成立. 显然,\(\bigcup\limits_{z \in K} B\left(z, \frac{\delta}{2}\right)\) 是 \(K\) 的一个开覆盖,根据有限覆盖定理,可以选出有限个圆盘,设为 \(B\left(z_j, \frac{\delta}{2}\right), j=1, \cdots, n\) 来覆盖 \(K\). 今取 \(\zeta_j \in B\left(z_j, \frac{\delta}{2}\right), j=1, \cdots, n\),由于 \(\{f_{n_k}(\zeta_j)\}, j=1, \cdots, n\) 收敛,由 Cauchy 收敛原理,存在 \(N\),当 \(l, k > N\) 时,有
\begin{align}
|f_{n_l}(\zeta_j) - f_{n_k}(\zeta_j)| < \varepsilon, \quad j=1, \cdots, n. \label{eq:cauchy_zeta}
\end{align}
今任取 \(z \in K\),\(z\) 必定属于 \(B\left(z_j, \frac{\delta}{2}\right), j=1, \cdots, n\) 中的一个,不妨设 \(z \in B\left(z_p, \frac{\delta}{2}\right)\). 因为 \(\zeta_p \in B\left(z_p, \frac{\delta}{2}\right)\),于是 \(|z - \zeta_p| < \delta\),因而
\begin{align}
|f_n(z) - f_n(\zeta_p)| < \varepsilon, \quad n=1, 2, \cdots. \label{eq:equi_cont}
\end{align}
由 \eqref{eq:cauchy_zeta} 式和 \eqref{eq:equi_cont} 式即知,当 \(l, k > N\) 时,便有
\begin{align*}
|f_{n_l}(z) - f_{n_k}(z)| \leqslant |f_{n_l}(z) - f_{n_l}(\zeta_p)| + |f_{n_l}(\zeta_p) - f_{n_k}(\zeta_p)| \quad + |f_{n_k}(\zeta_p) - f_{n_k}(z)| < 3\varepsilon.
\end{align*}
这就证明了 \(\{f_{n_k}\}\) 在 \(K\) 上一致收敛.
\end{proof}

\begin{theorem}[Montel定理]\label{theorem:复变-----Montel定理}
设 \(\mathscr{F}\) 是域 \(D\) 上的全纯函数族,那么 \(\mathscr{F}\) 是正规族的充分必要条件是 \(\mathscr{F}\) 在 \(D\) 上内闭一致有界.
\end{theorem}
\begin{proof}

{\heiti 必要性:}如果 \(\mathscr{F}\) 是 \(D\) 上的正规族,但不是内闭一致有界的,那么存在一个紧集 \(K \Subset D\),使得 \(\sup \{|f(z)|: z \in K, f \in \mathscr{F}\} = \infty.\) 因而存在序列 \(\{f_n\} \subseteq \mathscr{F}\),使得
\begin{align}
\sup \{|f_n(z)|: z \in K\} \geqslant n. \label{eq:sup_n}
\end{align}
由于 \(\mathscr{F}\) 是 \(D\) 上的正规族,故在 \(\{f_n\}\) 中存在子列 \(\{f_{n_k}\}\),它在 \(D\) 上内闭一致收敛,设其极限函数为 \(f\),则当 \(k > k_0\) 时,\(|f_{n_k}(z) - f(z)| < 1\) 在 \(K\) 上成立. \(f\) 是 \(D\) 上的全纯函数,当然在 \(K\) 上有界,不妨设 \(|f(z)| \leqslant M \quad (z \in K)\). 于是,当 \(z \in K\) 且 \(k > k_0\) 时,有
\begin{align}
|f_{n_k}(z)| \leqslant |f_{n_k}(z) - f(z)| + |f(z)| < M + 1. \label{eq:bound_n}
\end{align}
由 \eqref{eq:sup_n} 式和 \eqref{eq:bound_n} 式就得到 \(n_k \leqslant M + 1\) 的矛盾!

{\heiti 充分性:}任取 \(\{f_n\} \subseteq \mathscr{F}\),再取圆盘 \(\overline{B(z_0, r)} \Subset D\),按假定,\(\{f_n\}\) 在 \(\overline{B(z_0, r)}\) 上一致有界. 利用引理 4.1.8,即知 \(\{f_n'\}\) 在 \(\overline{B(z_0, r)}\) 上也一致有界,不妨设 \(|f_n'(z)| \leqslant M, \quad z \in \overline{B(z_0, r)}, n=1, 2, \cdots.\) 对任意 \(\varepsilon > 0\),取 \(\delta = \frac{\varepsilon}{M}\),当 \(z_1, z_2 \in B(z_0, r)\),且 \(|z_1 - z_2| < \delta\) 时,有
\begin{align*}
|f_n(z_1) - f_n(z_2)| = \left| \int_{z_2}^{z_1} f_n'(z) \mathrm{d}z \right| \leqslant M |z_1 - z_2| < \varepsilon.
\end{align*}
这就证明了 \(\{f_n\}\) 在 \(B(z_0, r)\) 上等度连续. 今取任意紧集 \(K \Subset D\),由于 \(K\) 可用有限个上面这种圆盘来覆盖,因而 \(\{f_n\}\) 在 \(K\) 上也是等度连续的,当然 \(\{f_n\}\) 在 \(K\) 上也是一致有界的. 由定理 \eqref{thm:arzela_ascoli},\(\{f_n\}\) 有子列 \(\{f_{n_k}\}\) 在 \(K\) 上一致收敛. 为了取出在任意紧集上一致收敛的子列,取一列紧集 \(K_j\),使得
\begin{align*}
K_1 \Subset K_2 \Subset \cdots \Subset K_j \Subset \cdots \to D.
\end{align*}
根据上面的讨论,先在 \(\{f_n\}\) 中取出在 \(K_1\) 上一致收敛的子列 \(\{f_n^{(1)}\}\),再从 \(\{f_n^{(1)}\}\) 中取出在 \(K_2\) 上一致收敛的子列 \(\{f_n^{(2)}\}\),继续这样做下去. 最后,用对角线法即可取出在所有 \(K_j\) 上一致收敛的子列,这个子列便在任意紧集上一致收敛.
\end{proof}



\end{document}